\section{Комбинаторика}
\label{sec:combinatorics}

\subsection{Правила суммы, произведения и биекции}

Комбинаторика изучает способы подсчёта конечных дискретных конфигураций. Три
основных приёма:
\begin{itemize}
 \item \textbf{правило суммы}: если классы объектов попарно не пересекаются и
       имеют мощности $N_1,\ldots,N_k$, то мощность их объединения равна
       $N_1+\cdots+N_k$;
 \item \textbf{правило произведения}: если объект строится последовательностью
       шагов и на $i$-м шаге всегда имеется $N_i$ вариантов, то всего вариантов
       $N_1\cdots N_k$;
 \item \textbf{принцип биекции}: если между двумя классами объектов построена
       биекция, то их мощности одинаковы.
\end{itemize}

\textbf{Принцип Дирихле.} Если $N$ объектов распределены по $m$ ящикам, то
какой-либо ящик содержит не менее
\[
    \left\lceil\frac Nm\right\rceil
\]
объектов. Простейшая форма: если объектов больше, чем ящиков, два объекта
окажутся в одном ящике.

\subsection{Перестановки, размещения и сочетания}

Для $n$ различных элементов:
\[
    P_n=n!.
\]
Если имеются группы одинаковых элементов размеров $n_1,\ldots,n_s$,
$n_1+\cdots+n_s=n$, то
\[
    P(n_1,\ldots,n_s)=\frac{n!}{n_1!\cdots n_s!}.
\]

Упорядоченный выбор $k$ различных элементов из $n$:
\[
    A_n^k=n(n-1)\cdots(n-k+1)=\frac{n!}{(n-k)!}.
\]
Если повторения разрешены, имеется $n^k$ последовательностей длины $k$.

Неупорядоченный выбор без повторений:
\[
    C_n^k=\binom nk=\frac{n!}{k!(n-k)!}.
\]
Связь
\[
    A_n^k=k!\binom nk
\]
объясняется тем, что каждый выбранный $k$-элементный набор можно упорядочить
$k!$ способами.

\subsection{Сочетания с повторениями. Метод ``звёзд и перегородок''}

Число решений в неотрицательных целых числах уравнения
\[
    x_1+\cdots+x_n=k,
    \qquad x_i\ge0,
\]
равно
\[
    \binom{n+k-1}{k}.
\]
Доказательство: расположим $k$ одинаковых ``звёзд'' и $n-1$ перегородок;
положение перегородок однозначно определяет числа $x_i$.

Если требуется $x_i\ge1$, делаем замену $y_i=x_i-1$ и получаем число решений
\[
    \binom{k-1}{n-1},\qquad k\ge n.
\]
Эта техника покрывает большое число задач о распределении одинаковых объектов.

\subsection{Биномиальные коэффициенты}

Бином Ньютона:
\[
 (x+y)^n=\sum_{k=0}^{n}\binom nk x^k y^{n-k}.
\]
Основные тождества:
\[
 \binom nk=\binom n{n-k},
 \qquad
 \binom nk=\binom{n-1}{k-1}+\binom{n-1}{k}.
\]
Последнее тождество (Паскаля) имеет простой комбинаторный смысл: при выборе
$k$ элементов из $n$ фиксированный элемент либо выбран, либо нет.

Подставляя $x=y=1$, получаем
\[
    \sum_{k=0}^n\binom nk=2^n,
\]
что одновременно является подсчётом всех подмножеств $n$-элементного множества.
При $x=1$, $y=-1$:
\[
    \sum_{k=0}^n(-1)^k\binom nk=0,
    \qquad n\ge1.
\]

Мультиномиальная формула:
\[
 (x_1+\cdots+x_s)^n=
 \sum_{n_1+\cdots+n_s=n}
 \frac{n!}{n_1!\cdots n_s!}
 x_1^{n_1}\cdots x_s^{n_s}.
\]

\subsection{Формула включений--исключений}

Для конечных множеств $A_1,\ldots,A_m$
\[
\left|\bigcup_{i=1}^{m}A_i\right|
=
\sum_i|A_i|
-\sum_{i<j}|A_i\cap A_j|
+\sum_{i<j<k}|A_i\cap A_j\cap A_k|
-\cdots
+(-1)^{m+1}|A_1\cap\cdots\cap A_m|.
\]

Почему формула работает: элемент, лежащий ровно в $r$ множествах, учитывается
с коэффициентом
\[
 \binom r1-\binom r2+\cdots+(-1)^{r+1}\binom rr=1.
\]

Классический пример --- число беспорядков, то есть перестановок без неподвижных
точек. Пусть $A_i$ --- событие ``элемент $i$ остался на месте''. Тогда
\[
    D_n=n!\sum_{k=0}^{n}\frac{(-1)^k}{k!},
\]
и асимптотически
\[
    D_n\sim \frac{n!}{e}.
\]

\subsection{Разбиения, числа Стирлинга и Белла}

Число разбиений $n$-элементного множества на $k$ непустых неупорядоченных
блоков --- число Стирлинга второго рода $S(n,k)$. Для последнего элемента есть
два случая: он либо образует новый блок, либо включается в один из уже имеющихся
$k$ блоков. Поэтому
\[
    S(n,k)=S(n-1,k-1)+kS(n-1,k).
\]
Граничные значения:
\[
    S(n,1)=S(n,n)=1.
\]

Явная формула, получаемая включениями--исключениями:
\[
    S(n,k)=\frac1{k!}\sum_{j=0}^{k}(-1)^{k-j}\binom kj j^n.
\]
Здесь $k!S(n,k)$ --- число сюръекций из $n$-элементного множества на
$k$-элементное.

Число всех разбиений --- число Белла
\[
    B_n=\sum_{k=0}^{n}S(n,k).
\]

Беззнаковые числа Стирлинга первого рода подсчитывают перестановки $n$ элементов
с заданным числом циклов.

\subsection{Производящие функции}

Обычная производящая функция последовательности $a_n$:
\[
    A(z)=\sum_{n\ge0}a_nz^n.
\]
Умножение производящих функций соответствует свёртке коэффициентов:
если $C(z)=A(z)B(z)$, то
\[
    c_n=\sum_{k=0}^{n}a_k b_{n-k}.
\]
Именно поэтому производящие функции удобны для подсчёта объектов, составленных
из независимых частей.

Для Фибоначчи $F_0=0$, $F_1=1$, $F_{n+2}=F_{n+1}+F_n$:
\[
 F(z)=\sum_{n\ge0}F_nz^n
 =\frac{z}{1-z-z^2}.
\]
Разложение рациональной функции на простейшие дроби приводит к формуле Бине.

Экспоненциальная производящая функция
\[
    \widehat A(z)=\sum_{n\ge0}a_n\frac{z^n}{n!}
\]
особенно естественна для помеченных комбинаторных объектов.

\subsection{Линейные рекуррентные соотношения}

Для
\[
 a_n=c_1a_{n-1}+\cdots+c_pa_{n-p}
\]
строится характеристический многочлен
\[
    \lambda^p-c_1\lambda^{p-1}-\cdots-c_p=0.
\]
Если корни $\lambda_i$ простые,
\[
    a_n=\sum_i C_i\lambda_i^n.
\]
Корню кратности $r$ соответствуют слагаемые
\[
    (C_0+C_1n+\cdots+C_{r-1}n^{r-1})\lambda^n.
\]
Константы определяются начальными условиями.

\subsection{Что важно сказать на экзамене}

В начале нужно классифицировать выбор: важен ли порядок, разрешены ли повторения.
Основные формулы --- $n!$, $A_n^k$, $\binom nk$, сочетания с повторениями и
включения--исключения. Для расширенного ответа полезны принцип Дирихле, числа
Стирлинга/Белла, производящие функции и характеристический метод решения
рекуррентных соотношений.

\newpage
